\section{Топологические группы}

\begin{definition}[7.1]
Множество $G$ со структурой группы и топологией называется \emph{топологической группой}, если отображения
\[
(g,h)\mapsto gh\colon G\times G\to G,\qquad g\mapsto g^{-1}\colon G\to G
\]
непрерывны (на $G\times G$~--- топология произведения).
\end{definition}

Пусть $a\in G$. Левый сдвиг $a_L\colon g\mapsto ag$ непрерывен как композиция непрерывных отображений
\[
\begin{tikzcd}[column sep=large]
G\arrow[r,"{g\,\mapsto\,(a,g)}"] & G\times G\arrow[r,"{(g,h)\,\mapsto\,gh}"] & G,
\end{tikzcd}
\]
и у него непрерывное обратное $(a^{-1})_L$. Значит, $a_L$~--- \emph{гомеоморфизм}; так же
гомеоморфизмами являются правые сдвиги $g\mapsto ga$, а также обращение $g\mapsto g^{-1}$. Отсюда:

\begin{lemma}[7.1a]
Пусть $H$~--- подгруппа топологической группы $G$.
\begin{enumerate}[label=\textup{(\alph*)}]
\item Смежный класс $aH$ открыт (замкнут), если $H$ открыта (замкнута).
\item Открытая подгруппа замкнута.
\item Замкнутая подгруппа конечного индекса открыта.
\item Если $G$ компактна, то открытая подгруппа имеет конечный индекс.
\item Подгруппа, содержащая открытую подгруппу (или вообще открытую окрестность единицы), открыта.
\end{enumerate}
\end{lemma}

\begin{proof}
(a) $aH=a_L(H)$, а $a_L$~--- гомеоморфизм.
(b) $G\setminus H=\bigcup_{a\notin H}aH$~--- объединение открытых множеств.
(c) $G\setminus H$~--- \emph{конечное} объединение замкнутых смежных классов $aH$, $a\notin H$, т.\,е.\ замкнуто.
(d) Смежные классы $aH$ открыты, попарно не пересекаются и покрывают $G$. Из этого покрытия можно выбрать конечное
подпокрытие, но ни один класс нельзя выбросить; значит, их конечное число.
(e) Если $U\subset H$~--- открытая окрестность $e$, то $H=\bigcup_{h\in H}hU$ открыта.
\end{proof}

Напомним, что \emph{окрестностью} точки $x$ называется любое множество, содержащее открытое множество, содержащее $x$
(окрестность не обязана быть открытой), а \emph{база окрестностей} точки $x$~--- это семейство $\mathcal N$ окрестностей
$x$ такое, что любое открытое
множество, содержащее $x$, содержит некоторое $N\in\mathcal N$.

\begin{proposition}[7.2]
Пусть $G$~--- топологическая группа и $\mathcal N$~--- база окрестностей единицы $e$. Тогда\footnote{Для подмножеств
$S,S'\subset G$ полагаем $SS'=\{ss'\mid s\in S,\ s'\in S'\}$ и $S^{-1}=\{s^{-1}\mid s\in S\}$.}
\begin{enumerate}[label=\textup{(\alph*)}]
\item для любых $N_1,N_2\in\mathcal N$ найдётся $N'\in\mathcal N$ такое, что $e\in N'\subset N_1\cap N_2$;
\item для любого $N\in\mathcal N$ найдётся $N'\in\mathcal N$ такое, что $N'N'\subset N$;
\item для любого $N\in\mathcal N$ найдётся $N'\in\mathcal N$ такое, что $N'\subset N^{-1}$;
\item для любых $N\in\mathcal N$ и $g\in G$ найдётся $N'\in\mathcal N$ такое, что $N'\subset gNg^{-1}$;
\item для любого $g\in G$ семейство $\{gN\mid N\in\mathcal N\}$~--- база окрестностей точки $g$.
\end{enumerate}
Обратно, пусть $G$~--- группа и $\mathcal N$~--- непустое семейство подмножеств $G$, удовлетворяющее (a)--(d). Тогда на
$G$ существует единственная топология, для которой выполнено (e); в этой топологии $G$~--- топологическая группа, а
$\mathcal N$~--- база окрестностей единицы.
\end{proposition}

\begin{center}
\begin{tikzpicture}[scale=0.95,every node/.style={font=\small}]
\def\blobN{(1.2,0.15) (0.85,0.72) (0.05,0.85) (-0.62,0.55) (-0.88,-0.05) (-0.45,-0.62) (0.4,-0.66) (1.02,-0.4)}
\def\blobNinv{(-1.2,-0.15) (-0.85,-0.72) (-0.05,-0.85) (0.62,-0.55) (0.88,0.05) (0.45,0.62) (-0.4,0.66) (-1.02,0.4)}
\def\blobL{(1.5,0.05) (1.05,0.56) (0.0,0.63) (-1.05,0.52) (-1.5,-0.02) (-1.0,-0.56) (0.05,-0.62) (1.05,-0.52)}
\def\Nprime#1{\fill[cnew!22] (0,0) circle (#1); \draw[cnew] (0,0) circle (#1);
  \fill (0,0) circle (1.3pt) node[right=0pt,inner sep=0.8pt,font=\scriptsize]{$e$};}
% (a)
\begin{scope}
\draw[cold,thick] (-0.48,0) ellipse (1.25 and 0.85); \draw[cold,thick] (0.48,0) ellipse (1.25 and 0.85);
\Nprime{0.3}
\node[cold] at (-1.55,0.95) {$N_1$}; \node[cold] at (1.55,0.95) {$N_2$}; \node[cnew,font=\scriptsize] at (0,0.5) {$N'$};
\node at (0,-1.45) {(a) $e\in N'\subset N_1\cap N_2$};
\end{scope}
% (b)
\begin{scope}[xshift=5.1cm]
\begin{scope}[scale=1.22] \draw[cold,thick] plot[smooth cycle] coordinates {\blobN}; \end{scope}
\draw[cnew,dashed] plot[smooth cycle] coordinates {(0.8,0.1) (0.5,0.62) (-0.25,0.66) (-0.72,0.12) (-0.42,-0.55) (0.4,-0.58)};
\Nprime{0.28}
\node[cold] at (1.6,0.95) {$N$}; \node[cnew,font=\scriptsize] at (0.1,0.42) {$N'N'$};
\node[cnew,font=\scriptsize] at (-0.53,-0.05) {$N'$};
\node at (0,-1.45) {(b) $N'N'\subset N$};
\end{scope}
% (c)
\begin{scope}[xshift=10.2cm]
\draw[cold,thick] plot[smooth cycle] coordinates {\blobN};
\draw[cold,dashed] plot[smooth cycle] coordinates {\blobNinv};
\Nprime{0.28}
\node[cold] at (1.3,0.8) {$N$}; \node[cold] at (-1.4,-0.8) {$N^{-1}$};
\node[cnew,font=\scriptsize] at (0,0.45) {$N'$};
\node at (0,-1.45) {(c) $N'\subset N^{-1}$};
\end{scope}
% (d)
\begin{scope}[xshift=2.55cm,yshift=-3.6cm]
\draw[cold,thick] plot[smooth cycle] coordinates {\blobL};
\begin{scope}[rotate=60] \draw[cold,dashed] plot[smooth cycle] coordinates {\blobL}; \end{scope}
\Nprime{0.28}
\node[cold] at (1.8,0.45) {$N$}; \node[cold,anchor=west] at (0.95,1.32) {$gNg^{-1}$};
\node[cnew,font=\scriptsize] at (0,0.42) {$N'$};
\node at (0,-1.75) {(d) $N'\subset gNg^{-1}$};
\end{scope}
% (e)
\begin{scope}[xshift=7.35cm,yshift=-3.6cm]
\begin{scope}[scale=0.62] \draw[cold,thick] plot[smooth cycle] coordinates {\blobN}; \end{scope}
\fill (0,0) circle (1.3pt) node[below=0.5pt,font=\scriptsize]{$e$};
\node[cold] at (-0.75,0.6) {$N$};
\begin{scope}[shift={(2.6,0.45)},scale=0.62] \draw[cnew,thick] plot[smooth cycle] coordinates {\blobN}; \end{scope}
\fill (2.6,0.45) circle (1.3pt) node[below=0.5pt,font=\scriptsize]{$g$};
\node[cnew] at (3.55,1.0) {$gN$};
\draw[->,>=stealth,cnew] (0.3,0.32) to[bend left=30] node[above]{$g_L$} (2.15,0.85);
\node at (1.3,-1.75) {(e) $gN$~--- окрестность $g$};
\end{scope}
\end{tikzpicture}
\end{center}
\figcap{Условия (a)--(e) в картинках: окрестности единицы можно пересекать, <<делить пополам>> относительно умножения,
обращать, сопрягать и переносить сдвигами в любую точку. В (c) множество $N^{-1}$~--- <<отражение>> $N$; оно не может
лежать строго внутри $N$: обращение~--- биекция, сохраняющая включения, и из $N^{-1}\subset N$ следует $N\subset N^{-1}$,
т.\,е.\ $N^{-1}=N$.}

\begin{proof}
\emph{1) Прямое утверждение.} (a) следует из определений: $N_1\cap N_2$ содержит открытое множество, содержащее $e$, а
оно содержит некоторое $N'$. (b): умножение непрерывно в точке $(e,e)$, поэтому для $N$ найдётся открытая окрестность
$W_1\times W_2$ точки $(e,e)$ с $W_1W_2\subset N$; возьмём $N'\subset W_1\cap W_2$ (по (a) и определению базы), тогда
$N'N'\subset N$. (c) так же следует из непрерывности $g\mapsto g^{-1}$ в $e$, (d)~--- из непрерывности
$h\mapsto g^{-1}hg$ в $e$ (её образ в $e$ равен $e$, а $\{h: g^{-1}hg\in N\}=gNg^{-1}$). (e) следует из того, что $g_L$~---
гомеоморфизм, переводящий $e$ в $g$.

\emph{2) Обратное: топология.} Заметим, что $e\in N$ для каждого $N\in\mathcal N$ (применим (a) к $N_1=N_2=N$).
Пусть $\mathcal U$~--- семейство подмножеств $U\subset G$ таких, что
\[
\text{для каждого }g\in U\text{ найдётся }N\in\mathcal N\text{ с }gN\subset U.
\]
Ясно, что $\varnothing,G\in\mathcal U$ и что объединение множеств из $\mathcal U$ лежит в $\mathcal U$. Если
$U_1,U_2\in\mathcal U$ и $g\in U_1\cap U_2$, то $gN_1\subset U_1$, $gN_2\subset U_2$, а по (a) есть $N'\subset N_1\cap N_2$,
так что $gN'\subset U_1\cap U_2$. Значит, $\mathcal U$~--- топология.

\emph{3) Условие (e) выполнено.} Нужно проверить две вещи: (e1) каждое $gN$~--- окрестность $g$, т.\,е.\ содержит
открытое множество, содержащее $g$; (e2) каждое открытое $U\ni g$ содержит некоторое $gN$. Условие (e2)~--- это
определение $\mathcal U$. Для (e1) положим
\[
V=\{x\in G\mid xN'\subset gN\ \text{для некоторого }N'\in\mathcal N\}.
\]
Тогда $g\in V$ (возьмём $N'=N$); $V\subset gN$, так как $x=xe\in xN'$; и $V$ открыто: если $xN'\subset gN$, выберем по (b)
$N''$ с $N''N''\subset N'$; для $y\in xN''$ имеем $yN''\subset xN''N''\subset xN'\subset gN$, т.\,е.\ $y\in V$, и
значит $xN''\subset V$. Итак, $g\in V\subset gN$ и $V\in\mathcal U$.

\emph{4) Единственность.} В любом топологическом пространстве множество открыто тогда и только тогда, когда оно
является окрестностью каждой своей точки. Если топология удовлетворяет (e), то открытость $U$ означает в точности, что
для каждого $g\in U$ найдётся $N$ с $gN\subset U$, т.\,е.\ $U\in\mathcal U$. Значит, такая топология единственна.

\emph{5) Умножение непрерывно.} Множества $g_1N_1\times g_2N_2$ образуют базу окрестностей точки $(g_1,g_2)$ в
$G\times G$. Пусть $U$ открыто и $g_1g_2\in U$; найдём $N_1,N_2$ с $g_1N_1g_2N_2\subset U$. Возьмём $N$ с
$g_1g_2N\subset U$, затем по (b) $N'$ с $N'N'\subset N$, и по (d) (для элемента $g_2$) $N_1\subset g_2N'g_2^{-1}$. Тогда
\[
g_1N_1g_2N'\subset g_1(g_2N'g_2^{-1})g_2N'=g_1g_2N'N'\subset g_1g_2N\subset U .
\]

\emph{6) Взятие обратного непрерывно.} Пусть $U$ открыто и $g^{-1}\in U$; найдём $N''$ с $(gN'')^{-1}\subset U$. Возьмём
$N$ с $g^{-1}N\subset U$; по (d) (для элемента $g^{-1}$) найдётся $N_3\subset g^{-1}Ng$, а по (c)~--- $N''\subset N_3^{-1}$.
Тогда
\[
(gN'')^{-1}=N''^{-1}g^{-1}\subset N_3g^{-1}\subset g^{-1}Ng\,g^{-1}=g^{-1}N\subset U.
\]
Наконец, $\mathcal N$~--- база окрестностей $e$ по (e) при $g=e$.
\end{proof}

\begin{remark}[7.2a]
\begin{enumerate}[label=(\alph*),leftmargin=*]
\item Каждое условие отвечает за свою часть: (a)~--- за то, что получается топология; (b)~--- за непрерывность умножения
в единице (и за то, что сдвиги $gN$ являются окрестностями); (c)~--- за непрерывность обратного в единице; (d)~--- за
перенос окрестностей из единицы в любую точку (непрерывность сопряжения). Без (b) утверждение неверно: для $G=\Z$ (аддитивно) и
$\mathcal N=\{\{-1,0,1\}\}$ условия (a), (c), (d) выполнены, а (b)~--- нет: $\{-1,0,1\}+\{-1,0,1\}=\{-2,\dots,2\}$.
Открытыми в $\mathcal U$ оказываются лишь $\varnothing$ и $\Z$ (если $U\ni g$ открыто, то $g\pm1\in U$, и по индукции
$U=\Z$), так что $\{g-1,g,g+1\}$ не является окрестностью $g$: топологии с условием (e) нет.
\item Полученная топология хаусдорфова тогда и только тогда, когда $\bigcap_{N\in\mathcal N}N=\{e\}$. Например,
семейство $\mathcal N=\{G\}$ удовлетворяет (a)--(d) и даёт антидискретную топологию.
\item Если все $N\in\mathcal N$~--- \emph{нормальные подгруппы}, то (b), (c), (d) выполняются с $N'=N$, и остаётся
проверить лишь (a). Именно так будет обстоять дело с семейством $\{G(S)\mid S\text{ инвариантно}\}$ в группе Галуа (\rf{7.6}).
\end{enumerate}
\end{remark}
